\section{Introduction}\label{sec:intro}
Link adaptation (LA) selects the Modulation and Coding Scheme (MCS) for downlink transmission from delayed and noisy Channel Quality Indicator (CQI) reports. In cellular networks (LTE and 5G New Radio), Outer Loop Link Adaptation (OLLA)~\cite{nakamura2002olla} serves as the ubiquitous industrial baseline, applying an additive offset that steps down on ACK and steps up on NACK to track a target Block Error Rate (BLER). In recent years, dozens of studies have proposed replacing or augmenting OLLA with reinforcement learning (RL) or bandit algorithms, reporting throughput and reliability improvements ranging from 10\,\% to over 100\,\%~\cite{saxena2022,ye2023,you2026,wang2026lolla,kela2022,demirel2025}.

Despite these promising reports, our screening of the literature reveals a critical comparator gap: of \NLitN{} learning-based link adaptation papers reviewed from available texts and abstracts, \NLitCalYes{} compared against a rule calibrated on the learner's own observation cues, \NLitCalPart{} did so partly, and for \NLitCalUnc{} the text did not permit a determination. In most cases, machine learning agents observing rich contextual information (such as Doppler velocity, delay history, or buffer status) were benchmarked solely against uncalibrated or global OLLA implementations. This leaves a fundamental question unresolved: \emph{does reinforcement learning outperform a deterministic heuristic given access to the same information?}

In this paper, we address this question through a pre-registered simulation campaign on an LTE/NR-like multi-user system with realistic multi-process HARQ feedback timing (RTT 8 TTIs). We evaluate learned policies against a comprehensive suite of calibrated deterministic baselines using strictly disjoint calibration, validation, and evaluation splits, paired common random numbers (CRN)~\cite{glasserman1992crn}, multiple training seeds, and a public case registry.

Our findings reveal both a clear positive domain and strict boundaries:
\begin{itemize}
\item \textbf{Contribution 1 (Methodology and Open Artefacts):} We formulate a pre-registered evaluation protocol for learning-based link adaptation with paired testing, multiplicity tracking, and an open case registry of \NRegCases{} cases and \NRegRows{} scheme evaluations.
\item \textbf{Contribution 2 (A Genuine Gain Under Realistic HARQ):} For deadline-constrained traffic where decisions depend dynamically on remaining packet slack and occupied HARQ retransmission buffers, a residual A2C policy \wordingsel{C2}{G6}{verdict}{unrun}.
\item \textbf{Contribution 3 (Where RL Does Not Help):} Under full-buffer goodput (M3, M4), reinforcement learning offers no advantage over calibrated static offsets ($-0.88$\,\% and $-1.73$\,\% vs.\ Fixed), with the learned policy collapsing to identity decisions.
\item \textbf{Contribution 4 (Pitfalls and Recommendations):} We document three simulation and framing errors that produced large illusory gains in earlier versions (mobility time-scale scaling, slot-stealing rewards, and PF-average timing artefacts under multi-process HARQ) and offer concrete reporting guidelines.
\end{itemize}
